 % Created by Robin Sinurat (Modified from LaTeX Arini)
% The University Center of Excellence for Advanced Intelligent Communications (AICOMS),
% Telkom University, Bandung, INDONESIA
% 26 March 2022

% Ukuran font 12pt
\documentclass[12pt, a4paper, onecolumn, oneside, final]{report}
% Load konfigurasi LaTeX untuk tipe laporan thesis
\usepackage{ta}
\usepackage{longtable}
\usepackage{tikz}
\usepackage{graphics}
\usepackage{subcaption}
\usepackage{times}
\usepackage{graphicx}
\usepackage[square,numbers,comma,sort,compress]{natbib}\usepackage{titlesec}
\usepackage{amsmath}
\usepackage{appendix}
\onehalfspacing
\usepackage{hyperref} 
\usepackage{xcolor,colortbl}
\usepackage{multicol}
\usepackage[graphicx]{realboxes}
\usepackage[nodayofweek]{datetime}
\usepackage{multirow}

\renewcommand{\contentsname}{CONTENTS}
\renewcommand{\chaptername}{CHAPTER}
\renewcommand{\listfigurename}{LIST OF FIGURES}
\renewcommand{\listtablename}{LIST OF TABLES}
\renewcommand{\appendixname}{APPENDIX}
\renewcommand{\bibname}{REFERENCES}
\renewcommand{\figurename}{Fig.}

\usepackage[labelsep=space]{caption} % to remove colon after figure
%\captionsetup{labelsep = period}
\captionsetup{labelsep=separatorbaru}
%\usepackage{array}
\newcolumntype{L}[1]{>{\raggedright\let\newline\\\arraybackslash\hspace{0pt}}m{#1}}
\newcolumntype{C}[1]{>{\centering\let\newline\\\arraybackslash\hspace{0pt}}m{#1}}
\newcolumntype{R}[1]{>{\raggedleft\let\newline\\\arraybackslash\hspace{0pt}}m{#1}} % to remove colon after figure

% Load konfigurasi khusus untuk laporan yang sedang dibuat
%-----------------------------------------------------------------------------%
% Informasi Mengenai Dokumen
%-----------------------------------------------------------------------------%
% 
% Judul laporan. 
\var{\judul}{\setstretch{1}Babibu}
% 
% Tulis kembali judul laporan, kali ini akan diubah menjadi huruf kapital
\Var{\Judul}{\setstretch{1}Tuliskan Judul Disini}

% Tipe laporan, dapat berisi Skripsi, Tugas Akhir, Thesis, atau Disertasi
\var{\type}{Master Thesis}
% 
% Tulis kembali tipe laporan, kali ini akan diubah menjadi huruf kapital
\Var{\Type}{Master Thesis}
% 

% Tulis nama penulis 
\Var{\Penulis}{Reynaldhi Tryana Graha}
% Tulis NPM penulis
\var{\nim}{201012420030}
\var{\alamat}{Rancasari, Bandung, Jawa Barat}
\var{\tlp}{+6285156964564}
\var{\email}{\textit{reynaldhitryanagraha@student.telkomuniversity.ac.id}}

% Tuliskan Fakultas dimana penulis berada
\Var{\Fakultas}{School of Electrical Engineering}
\var{\fakultas}{School of Electrical Engineering}
% 
% Tuliskan Program Studi yang diambil penulis
\Var{\Program}{Master of Electrical-Telecommunication Engineering}
\var{\program}{Master of Electrical-Telecommunication Engineering}
% 
% Tuliskan tahun publikasi laporan
\Var{\Tahun}{2024}
 
% Tuliskan tanggal pengesahan laporan, waktu dimana laporan diserahkan ke penguji/sekretariat
\var{\tanggalPengesahan}{\today} 

% Tuliskan pembimbing 
\var{\pembimbingSatu}{Dr. Koredianto Usman, S.T., M.Sc.}
\var{\nikSatu}{02750053}
\var{\pembimbingDua}{Suryo Adhi Wibowo, S.T., M.T., Ph.D.}
\var{\nikDua}{1087003}
% 
% Alias untuk memudahkan alur penulisan paa saat menulis laporan
\var{\saya}{Author}

%-----------------------------------------------------------------------------%
% Judul Setiap Bab
%-----------------------------------------------------------------------------%
% 
% Berikut ada judul-judul setiap bab. 
% Silahkan diubah sesuai dengan kebutuhan. 
% 
\var{\kataPengantar}{PREFACE}
\var{\babSatu}{INTRODUCTION}
\var{\babDua}{BASIC CONCEPT}
\var{\babTiga}{THE PROPOSED CODES}
\var{\babempat}{PERFORMANCE EVALUATIONS}
\var{\babLima}{CONCLUSION}


% Awal bagian penulisan laporan
\begin{document}

% Sampul Laporan
\begin{titlepage}
    \begin{center}      
    	\vspace*{0.15cm}
        % judul thesis harus dalam 14pt Times New Roman
        \bo{\large\Judul} \\[0.5cm]
        \vspace*{0cm}    
        % harus dalam 14pt Times New Roman
        \textbf{A MASTER'S THESIS}\\
        \vspace*{0.5cm}
		\setstretch{1}
		Submitted to\\
		Graduate School of Electrical Engineering
        \vspace*{1.5cm}
                      
        \begin{figure}
            \begin{center}
                \includegraphics[scale=0.35]{logo/logo_telu1}
            \end{center}
        \end{figure}
        \vspace*{0.75cm}
           % penulis dan npm
        By\\
        \vspace*{0.25cm}
        \bo{\Penulis} \\
        \vspace*{0.25cm}
        \bo{\nim} \\
        
        \vspace*{1.5cm}
        
        \setstretch{1}
        In partial fulfillment of the requirements\\
        for the Degree of Master of Engineering\\
        \vspace*{1.5cm}
        % informasi mengenai fakultas dan program studi
        \bo{
        	TELKOM UNIVERSITY\\
        	BANDUNG \\
        	\Tahun
        }
    \end{center}
\end{titlepage}


% Tidak mengaktifkan penomoran halaman
\pagenumbering{gobble}

% Halaman pengesahan
\addChapter{APPROVAL PAGE}
%-----------------------------------------------------------------------------
% Halaman Pengesahan
%-----------------------------------------------------------------------------


\begin{center}
	
	\textbf{\large{APPROVAL PAGE}}
	\vspace*{1 cm}
	
    \textbf{\large{MASTER'S THESIS}}
	\vspace*{1 cm}

    \textbf{\Judul}
%	\vspace*{1 cm}    
%    
%    \textit{\textbf{\JudulInggris}}
	\vspace*{1.5 cm}
    
    by\\
    \vspace*{0.25 cm}
    
    \textbf{\Penulis}\\
    \textbf{\nim}\\
   	\vspace*{2 cm}   
    
     {\bfseries Approved and authorized to fulfil one of the requirements of\\Program of \program\\ \fakultas\\ Telkom University\\Bandung
    }
 %    \\
	% \textbf{Program of \program}\\
 %    \textbf{\fakultas}\\
 %    \textbf{Telkom University}\\
 %    \textbf{Bandung}
 
    
    
    \vspace*{3 cm}
    
    \textbf{Bandung, \today}\\
%    \textbf{Menyetujui,}
    
\end{center}
    
\begin{tabular}
	{>{\centering\arraybackslash} p{0.3\paperwidth} >{\centering\arraybackslash} p{0.34\paperwidth}}\\
    Supervisor & Co-Supervisor\\[3 cm]
    \uline{\textbf{\pembimbingSatu}} & \uline{\textbf{\pembimbingDua}}\\
    NIP. \nikSatu & NIP. \nikDua
\end{tabular}

%-----------------------------------------------------------------------------

% Halaman pernyataan orisinalitas
\addChapter{SELF DECLARATION AGAINST PLAGIARISM}
\chapter*{}

%\AddToShipoutPicture*{%
%\gradientbox{white}{white}{%
%
%\begin{minipage}[ct][0.98\paperheight][t]{\paperwidth}%
%\end{minipage}}}
%
%    \begin{center}
%    \textbf{SELF DECLARATION AGAINST PLAGIARISM}\\
%    \end{center}
%    
%    \begin{tabular}{ll}
%    Nama & :\hspace*{0.2 cm}\penulis \\
%    NIM & :\hspace*{0.2 cm}\nim \\
%    Alamat & :\hspace*{0.2 cm}\alamat \\
%    No. Telepon & :\hspace*{0.2 cm}\tlp \\
%    Email & :\hspace*{0.2 cm}\email \\
%    \end{tabular}
%    
%    \vspace*{1 cm}
%    Menyatakan bahwa Tugas Akhir ini merupakan karya orisinal saya sendiri, dengan judul:
%    
%    \begin{center}
%    \textbf{\Judul}\\
%%    \textit{\textbf{\JudulInggris}}\\
%    \end{center}
%    
%    Atas pernyataan ini, saya siap menanggung resiko\slash sanksi yang dijatuhkan kepada saya apabila kemudian ditemukan adanya pelanggaran terhadap kejujuran akademik atau etika keilmuan dalam karya ini, atau ditemukan bukti yang menunjukkan ketidakaslian karya ini.
%    
%    \vspace*{1 cm}
%    
%    \begin{tabular}{cl}
%    \multirow{6}{*}{\includegraphics[scale=0.03,origin=c]{pics/fotoa}\hspace{4 cm}}
%    & Bandung, \tanggalPengesahan \\
%    & \includegraphics[scale=1,origin=c]{pics/sign_Arini} \\
%    & \\
%    & \penulis \\
%    \cline{2-2}
%    &  \nim\\
%    \end{tabular}
    
    \begin{center}
	\textbf{\large{SELF DECLARATION AGAINST PLAGIARISM}}
	\vspace*{1 cm}
\end{center}
	I hereby declare that all information in this document has been obtained and presented in accordance with academic rules and ethical conduct. I also declare that, as required by these rules and conduct. I have full cited and referenced all materials and results that are not original to this work.\\
	
	\noindent\today\\
	\Penulis\\
% 	\includegraphics[scale=1,origin=c]{pics/sign_Arini}
	
	\noindent Signature:\rule{4cm}{0.4pt}


% Menggunakan penomoran halaman Romawi
\pagenumbering{roman}

% Setelah bagian ini, halaman dihitung sebagai halaman ke-2
\setcounter{page}{4}

% Abstrak dalam Bahasa Inggris
\addChapter{ABSTRACT}

\chapter*{abstract}
\vspace*{-0.3cm}

\indent \indent 
Quantum technology in recent decades is attracting great attentions, especially on the quantum computing, quantum communications, and quantum key distribution (QKD), which are candidates for the sixth generation of telecommunications (6G) in 2030 and beyond.

\vspace{0.2in}
\noindent\textbf{Keywords:} CSS, Decoherence, Depolarizing Channels, Red-Muller codes, QWER.


\newpage

% Ucapan Terima Kasih
\addChapter{ACKNOWLEDGEMENTS}
%-----------------------------------------------------------------------------%
\chapter*{ACKNOWLEDGEMENTS}
%-----------------------------------------------------------------------------%
This thesis is compiled with the effort, help, and support from all supporting elements. The author would like to express the deepest gratitude and thanks to:
\begin{enumerate}
\item Allah SWT, for all the love, guidance and forgiveness in every mistake that the author has ever done and Rasulullah SAW, as role model who inspire writer in living life and trying to be better.

\end{enumerate}
  
%\vspace*{0.1cm}
%\begin{flushright}
%Bandung, 02 Februari 2019\\[0.1cm]
%\vspace*{1cm}
%\penulis

%\end{flushright}

% Kata Pengantar
\addChapter{\kataPengantar}
%-----------------------------------------------------------------------------%
\chapter*{\kataPengantar}
%-----------------------------------------------------------------------------%
	Alhamdu lillahi rabbil 'alamin, praise to Allah, the most gracious, the most merciful, with the mercy and guidance, the author has successfully finished this thesis with the title of ~"\textbf{\Judul}". The author compiled this thesis to be filled in the graduation requirements in Program of Master of Electrical-Telecommunication Engineering, School of Electrical Engineering, Telkom University.
		
	Some parts of this thesis also have been submitted to The 9th Asia Pacific Conference on Wireless \& Mobile (APWIMOB) Conference 2024.
	
	The suggestions for improving this thesis are highly appreciated. Hopefully, this thesis is expected to be improved and provided contributions for the reader and Indonesia especially for education and research of telecommunication on the future.   

\vspace*{1 cm}
\begin{flushright}
Bandung, \tanggalPengesahan\\
\vspace*{2.75 cm}

\Penulis

\end{flushright}
%-----------------------------------------------------------------------------

% Daftar Isi
\tableofcontents
\clearpage

% Daftar Gambar
\listoffigures
\clearpage

% Daftar Tabel
\listoftables
\clearpage

% Daftar Lampiran
\addChapter{LIST OF ABBREVIATION}
%-----------------------------------------------------------------------------%
\chapter*{LIST OF ABBREVIATION}
%-----------------------------------------------------------------------------%
\begin{tabular}{lll}
1G &:& First Generation \\
2G &:& Second Generation \\
3G &:& Third Generation \\
4G &:& Fourth Generation \\
5G &:& Fifth Generation \\
6G &:& Sixth Generation \\
CSS &:& Calderbank-Shor-Steane\\
MDS &:& Maximum distance separable\\
LUT &:& Look Up Table\\
QECC &:& Quantum Error Correction Codes\\
QKD &:& Quantum Key Distribution\\
QWER &:& Quantum Word Error Rate\\
RM &:& Reed-Muller codes\\
SIP &:& Simplectectic Inner Product\\
WP &:& Work Packages\\

\end{tabular}

% Daftar Lampiran
\addChapter{LIST OF SYMBOL}
%-----------------------------------------------------------------------------%
\chapter*{LIST OF SYMBOL}
%-----------------------------------------------------------------------------%
\begin{longtable}{lll}
$\alpha$            & Constant number complex\\
$\beta$             & Constant number complex\\
$\rho$              & The density operator\\
$\rho^\prime$       & Evolution of quantum system\\

\end{longtable}


% Daftar Achievement
\addChapter{ACHIEVEMENT}
%-----------------------------------------------------------------------------
% Daftar Achievement
%-----------------------------------------------------------------------------


\chapter*{Achievements}

\vspace*{0.5 cm}

\begin{enumerate}
\item AAAAA
\end{enumerate}

%-----------------------------------------------------------------------------

% Gunakan penomeran Arab (1, 2, 3, ...) setelah bagian ini.
\pagenumbering{arabic}

% Bab 1 : Pendahuluan
%-----------------------------------------------------------------------------%
\chapter{\babSatu}
%-----------------------------------------------------------------------------
\vspace*{0.2cm}
\hyphenation{se-ve-ral}

\section{Background}
%Background merisikan judul, motivasi dan masalah serta yang diusulkan
%dari paper
With information theory led the mathematician Peter Shor to introduce a fast-factoring algorithm for a quantum computing, a quantum computing whose bits exist in superposition and can be entangled \cite{symetri}. Superposition and entanglement are useful for quantum computing and quantum key distribution (QKD).
%
% Bab 2 : Dasar Teori
%----------------------------------------------------------------------
\chapter{\babDua}
%--------------------------------------------------------
This chapter presents some basic theory that serves as a foundation for proposes a new quantum coding scheme constructed from classical codes. 
%--------------------------------------------------------------------
\section{General Communications System}

\textcolor{black}{Communication systems is a system that allows the information to occur between two points. The main elements in the communication systems process are transmitters, channels, and receiver information}. Transmitter is one of the main elements in the communications process because at this stage settings are made to process data or signals into a form that is in accordance with the transmission. Generally, at this stage there are other stages of processing i.e., the encoding process, the modulation process, etc. 

%
% Bab 3 : Perancangan
%-------------------------------------------------------------------------
\chapter{\babTiga}
%--------------------------------------------------------------------------
This thesis proposes new quantum coding scheme 

% Bab 4 : Pengujian
\chapter{\babempat}
This section evaluates the performance of the proposed codes in several aspects (i) syndrome extraction produces unique codes for each single errors occurred in quantum and (ii) the performance of quantum word error rate (QWER). 

% Bab 5 : Kesimpulan dan Saran
%---------------------------------------------------------------
\chapter{\babLima}
%---------------------------------------------------------------
\section{Conclusions}
This thesis has proposed a new quantum coding scheme 

\section{Future Works}
{\color{black}This thesis can be further possible be developed, especially, in terms of QWER performance under the other quantum channels.}


%
\bibliographystyle{IEEEtran}
\bibliography{ref_database}
%\newpage
% Lampiran 
%\begin{appendix}
%	\pagenumbering{arabic}
%	%
% @author  Andreas Febrian
% @version 1.00 
% 
% Hanya sebuah pembatas bertuliskan LAMPIRAN ditengah halaman. 
% 

\begin{titlepage}
	\centering 
	\vspace*{6cm}
	\noindent \Huge{LAMPIRAN}
	\addChapter{LAMPIRAN}
\end{titlepage}

%%	\input{appendices_IEEE}
%	\input{appendices}
%	%\setcounter{page}{2}
%\end{appendix}
\end{document}